\section{Experiments}
\label{sec:experiments}

We evaluate \cls{}, \cwm{}, and trajectory-level scoring (TRAJ) across model sizes and complexity levels.

\subsection{Experimental Setup}
\label{sec:setup}

\paragraph{Datasets.}
\emph{Function-level}: mutated variants from HumanEval~\cite{chen2021humaneval} and MBPP~\cite{austin2021mbpp} (five mutation types), labels verified by execution. 19,076 samples (train 15,290 / val 1,896 / test 1,890; pass:fail $\approx$ 52:48, majority baseline 51.75\%).
\emph{Repo-level}: per-test instances from SWE-bench Verified~\cite{chowdhury2024swebenchverified} mutants. ${\sim}$70K training samples; val ($n{=}9{,}335$) and test ($n{=}9{,}784$) have zero problem overlap (majority baseline 72.96\%). Val-based checkpoint selection follows standard practice~\cite{jimenez2024swebench}.

\paragraph{Models and training.}
We fine-tune Qwen3-4B and Qwen3-8B~\cite{qwen2025qwen3} with LoRA~\cite{hu2022lora} (rank 16, 3 epochs) across five seeds. Seeds 42/123/456 use bfloat16; seeds 789/0 use NF4 quantization~\cite{dettmers2024qlora}. Our \cwm{} baseline is a LoRA fine-tuned variant for controlled comparison, not a reproduction of the 32B Code World Model~\cite{meta2025cwm}. Repo-level models use val-selected checkpoints. Zero-shot baselines use Qwen3-8B and Qwen3-32B (NF4) as LLM-as-judge.

\paragraph{Statistics.}
Confidence intervals use clustered bootstrap~\cite{davison1997bootstrap} (114 problem clusters, 10K resamples); per-test vs.\ trajectory-level comparison uses TOST~\cite{schuirmann1987tost} ($\pm$2\pp{} margin). ICC~\cite{shrout1979icc} quantifies within-problem correlation. Oracle analysis uses SWE-bench Verified (500 instances).


\subsection{Function-Level Results}
\label{sec:func_results}

Table~\ref{tab:main_results} reports function-level results across formulations.

\begin{table}[t]
\centering
\small
\caption{Function-level per-test verification results on HumanEval+MBPP test set (1,890 samples). \cls{}: per-test binary classification; \cwm{}: execution output generation; TRAJ: trajectory-level scoring; ZS: zero-shot LLM-as-judge. Mean $\pm$ std over seeds reported where applicable.}
\label{tab:main_results}
\resizebox{\columnwidth}{!}{%
\setlength{\tabcolsep}{4pt}
\begin{tabular}{llcc}
\toprule
\textbf{Method} & \textbf{Model} & \textbf{Seeds} & \textbf{Accuracy (\%)} \\
\midrule
\cls{} (fine-tuned) & 4B & 5 & \textbf{88.43 $\pm$ 0.78} \\
\cwm{} (fine-tuned) & 4B & 5 & 87.19 $\pm$ 0.37 \\
TRAJ (fine-tuned) & 4B & 5 & 86.30 $\pm$ 0.97 \\
\midrule
ZS LLM-as-judge & 8B & 1 & 64.40 \\
ZS LLM-as-judge & 32B (NF4) & 1 & 71.69 \\
\bottomrule
\end{tabular}}
\end{table}

Fine-tuned \cls{} achieves 88.43\% ($\pm$0.78), compared to 87.19\% ($\pm$0.37) for \cwm{} and 86.30\% ($\pm$0.97) for TRAJ. The \cls{}--\cwm{} gap is 1.24\pp{}, modest at this complexity level where generation is not yet bottlenecked by token budget or string complexity. Zero-shot baselines lag substantially (64.40\% at 8B, 71.69\% at 32B NF4), confirming that task-specific fine-tuning is necessary for reliable verification.


\subsection{Scaling Behavior}
\label{sec:scaling}

Table~\ref{tab:scaling} compares \cls{} and \cwm{} at two model sizes using three matched seeds (42, 123, 456), all trained in full bfloat16 precision to ensure a controlled comparison.

\begin{table}[t]
\centering
\small
\caption{Model-size scaling: \cls{} vs.\ \cwm{} accuracy (\%) at 4B and 8B. Three matched seeds in full precision. The \cls{}--\cwm{} gap widens from 1.18\pp{} to 2.43\pp{}.}
\label{tab:scaling}
\setlength{\tabcolsep}{4pt}
\begin{tabular}{llccccc}
\toprule
\textbf{Method} & \textbf{Size} & \textbf{s42} & \textbf{s123} & \textbf{s456} & \textbf{Mean} & \textbf{Gap} \\
\midrule
\cls{} & 4B & 87.83 & 89.52 & 87.88 & 88.41 & \multirow{2}{*}{1.18} \\
\cwm{} & 4B & 87.51 & 87.20 & 86.98 & 87.23 & \\
\midrule
\cls{} & 8B & \textbf{90.48} & 89.40 & 88.31 & \textbf{89.40} & \multirow{2}{*}{\textbf{2.43}} \\
\cwm{} & 8B & 87.67 & 86.98 & 86.24 & 86.96 & \\
\bottomrule
\end{tabular}
\end{table}

At 4B, \cls{} and \cwm{} are separated by 1.18\pp{} on the matched three-seed subset. At 8B, the gap more than doubles to 2.43\pp{} (Figure~\ref{fig:scaling_gap}). Clustered bootstrap analysis yields a mean gap of 2.26\pp{} ($p = 0.0066$, 95\% CI $[1.80, 2.71]$), confirming the widening is not a sampling artifact.

The asymmetry in scaling efficiency is striking: from 4B to 8B, \cls{} gains +0.99\pp{} (88.41 $\to$ 89.40) while \cwm{} loses $-$0.27\pp{} (87.23 $\to$ 86.96). The additional model capacity benefits classification---where the model needs only to discriminate pass from fail---but provides no measurable improvement for generation, where the model must additionally produce exact output strings. This pattern is consistent with the hypothesis that the string-prediction bottleneck, not model capacity, bounds \cwm{} accuracy. However, the \cwm{} regression ($-$0.27\pp{}) falls within seed noise (3-seed std ${\sim}$0.5\pp{}) and could reflect suboptimal hyperparameter transfer rather than a fundamental bottleneck. Whether this divergence amplifies at repository-level complexity is examined in \S\ref{sec:repo}.


\subsection{Per-Test vs.\ Trajectory-Level Equivalence}
\label{sec:equivalence}

We compare \cls{} (aggregated via all-pass rule) against TRAJ on 523 variant-level instances across five seeds. The pooled gap is $-$0.45\pp{} (95\% CI [$-$1.71, +0.84]); the TOST equivalence test confirms statistical equivalence at $p = 0.007$ ($\pm$2\pp{} margin). Within-problem predictions are correlated (ICC $= 0.191$, $p < 0.0001$, DEFF $= 3.97$); all statistical tests use clustered bootstrap accordingly. Per-test classification thus provides diagnostic information---which specific tests fail---at no cost to variant-level accuracy.


\subsection{Accuracy-Utility Frontier}
\label{sec:oracle-results}

Oracle analysis on SWE-bench Verified (500 instances) reveals that per-test advantage concentrates sharply: at $T = 1$ F2P tests (68.7\% of instances), per-test and trajectory-level are informationally equivalent. The advantage materializes at $T \geq 2$ (31.3\%), with $T = 2$ alone contributing 55.3\% of the total 10.93\pp{} block-correlated advantage. Full stratification and the i.i.d.\ vs.\ block-correlated comparison appear in Appendix~\ref{app:oracle}.


\subsection{Repo-Level Results}
\label{sec:repo}

Table~\ref{tab:repo_level} compares \cls{} and \cwm{} on SWE-bench Verified per-test classification at repository-level complexity. All checkpoints are selected by validation accuracy; the test set ($n{=}9{,}784$) has zero problem overlap with validation.

\begin{table}[t]
\centering
\small
\caption{Repo-level per-test verification on SWE-bench Verified test set ($n{=}9{,}784$, val-selected checkpoints). Same-size 4B comparison is primary; \cls{} 8B included for scaling context; \cwm{} 8B in progress.}
\label{tab:repo_level}
\resizebox{\columnwidth}{!}{%
\setlength{\tabcolsep}{3pt}
\begin{tabular}{llccccc}
\toprule
\textbf{Method} & \textbf{Size} & \textbf{Seed} & \textbf{Acc (\%)} & \textbf{Macro \fone{}} & \textbf{Unparse (\%)} & \textbf{vs Maj.\ (\pp{})} \\
\midrule
\cls{} & 4B & 42 & \textbf{83.63} & 75.63 & 0.0 & +10.67 \\
\cwm{} & 4B & 42 & 58.04 & 51.67 & 53.7 & $-$14.92 \\
\cwm{} & 4B & 123 & 59.09 & 43.43 & 67.2 & $-$13.87 \\
\midrule
\cls{} & 8B & 42 & 84.48 & 78.17 & 0.0 & +11.52 \\
\cls{} & 8B & 123 & 85.60 & 78.38 & 0.0 & +12.64 \\
\cwm{} & 8B & --- & --- & --- & --- & --- \\
\midrule
Majority & --- & --- & 72.96 & --- & --- & --- \\
\bottomrule
\end{tabular}}
\end{table}

The same-size 4B comparison reveals a 25.59\pp{} gap (\cls{} 83.63\% vs.\ \cwm{} 58.04\%, seed 42; Figure~\ref{fig:scaling_gap}). \cwm{} falls 14.92\pp{} below the majority baseline (72.96\%), indicating negative utility. A second \cwm{} seed (123, 59.09\%) confirms the failure is not seed-specific. The primary driver is format compliance collapse: 53.7--67.2\% of \cwm{} outputs are unparseable across seeds, while \cls{} produces zero unparseable outputs under standard inference (8,192 tokens). Among parseable \cwm{} outputs, accuracy is 42--43\%---below both majority baseline and chance level (50\%)---indicating anti-informative predictions. \S\ref{sec:analysis} provides detailed error analysis.

\cls{} 8B achieves 84.48\% (seed 42) and 85.60\% (seed 123) in cross-size comparison. The 4B$\to$8B transition within \cls{} (83.63\%$\to$84.48\%, seed 42) is within noise (+0.85\pp{}, single seed). \cwm{} 8B evaluation is in progress.

\paragraph{Sequence length.} A matched-context ablation confirms input context disparity explains $<$3\% of the \cls{}--\cwm{} gap: \cls{} 8B evaluated at \texttt{max\_seq\_length}$=$2,048 yields a matched-context gap of 25.00\pp{}, representing 98\% of the same-size 25.59\pp{} gap (Appendix~\ref{app:ablations}).

\paragraph{Limitations.} The same-size 4B comparison uses a single seed; \cls{} 4B seed 123 (val-best ckpt-800) achieves only 75.33\% test accuracy (+2.4\pp{} above majority baseline), indicating that val-best selection at an early checkpoint (step 800) amplifies the val-test distribution shift observed across all configurations. The validation and test sets have zero problem overlap but differ in mutation type composition, introducing potential distribution shift. All repo-level results use val-selected checkpoints, and checkpoint selection variance could contribute to observed differences. Our training data is synthetically constructed from mutation operators rather than real developer patches; the synthetic-to-real transfer gap remains untested, and classifier behavior on naturally occurring bugs may differ from the controlled mutation patterns studied here.


\subsection{Analysis}
\label{sec:analysis}

\paragraph{(a) Error complementarity.}
\cls{} and \cwm{} make errors in opposite directions at function level. Among \cls{} errors (seed 42), 75.5\% are false negatives (FAIL $\to$ PASS), meaning the classifier tends to over-predict passing. Among \cwm{} errors, 83.1\% are false positives (PASS $\to$ FAIL), reflecting string-matching failures that incorrectly reject correct outputs. An oracle combination that takes the correct answer whenever either model is right achieves 90.63\% accuracy, a 3.12\pp{} improvement over \cls{} alone (87.51\%). This complementarity suggests that ensemble approaches combining classification and generation signals could yield further gains.

\paragraph{(b) CWM error mode breakdown.}
Among unparseable \cwm{} outputs, 76.9\% contain no execution keywords; 22.9\% contain ambiguous co-occurring pass/fail keywords. Format compliance varies sharply across repositories (scikit-learn: near-zero unparseable; sphinx: 100\%), suggesting domain-specific code contexts prevent structured generation.

\paragraph{(c) Relaxed matching does not rescue \cwm{}.}
Relaxed matching (tolerating format deviations, L0--L4) \emph{decreases} \cwm{} accuracy: 58.04\%$\to$56.98\% (seed 42) and 59.09\%$\to$55.91\% (seed 123). Unparseable predictions default to the majority class (pass, 72.96\%), which outperforms the model's own parseable predictions (42--53\%). The \cls{}--\cwm{} gap widens from 26\pp{} to 28\pp{}, ruling out the hypothesis that parsing strictness inflates the observed gap.

\paragraph{(d) Variance compression.}
\cwm{}'s 5-seed standard deviation (0.37\pp{}) is lower than \cls{} (0.78\pp{}) and TRAJ (0.97\pp{}) at function level. This is consistent with the bottleneck hypothesis: if \cwm{} performance is capped by output-format constraints, seed-to-seed variability is compressed relative to formulations whose performance depends more on learned representations.

\paragraph{(e) Per-test precision and recall.}
\cls{} achieves fail \fone{} 87.86\% (precision 82.75\%, recall 93.64\%) at seed 42, with 42.1\% of problems perfectly predicted. Multi-F2P instances (6--10 tests) reach 96.3\% accuracy, corroborating the oracle finding.

\paragraph{(f) Block correlation.}
Within-problem errors are correlated (ICC $= 0.191$, DEFF $= 3.97$, 2.81$\times$ more all-correct problems than i.i.d.\ expectation). All tests use clustered bootstrap accordingly.

\paragraph{(g) LoRA rank and majority voting.}
Doubling adapter capacity (rank 32) does not improve \cwm{} accuracy ($-$0.51\pp{}, within seed variance); majority voting ($K{=}3$) decreases accuracy ($-$0.68\pp{}). Both confirm the bottleneck is formulation-level, not capacity-level (Appendix~\ref{app:ablations}).